%%%%%%%%%%%%%%%%%%%%%%%%%%%%%%%%%%%%%%%%%%%%%%%%%%%%%%%%%%%%%%%%%%%%%%%%
\begin{frame}[fragile]\frametitle{What is ROC?}

\begin{columns}
\begin{column}[T]{0.5\linewidth}
\begin{itemize}
\item ROC : Receiver Operating Characteristics Curve
\item ROC tells us about how good
the model can distinguish between two things (e.g If a patient has a
disease or no). 
\item Better models can accurately distinguish between the two.
\item Predictions are plotted on x, and colors show True or false. Y is frequency or counts.
\end{itemize}
\end{column}
\begin{column}[T]{0.5\linewidth}

\begin{center}
\includegraphics[width=0.9\linewidth]{roc1}
\end{center}

\tiny{(Ref: Let's learn about AUC ROC Curve! - GreyAtom)}
\end{column}

\end{columns}
\end{frame}

%%%%%%%%%%%%%%%%%%%%%%%%%%%%%%%%%%%%%%%%%%%%%%%%%%%%%%%%%%%%%%%%%%%%%%%%
\begin{frame}[fragile]\frametitle{ROC: Setting a Threshold}

\begin{columns}
\begin{column}[T]{0.5\linewidth}
\begin{itemize}
\item Pick a value where we need to set the cut off say, ``0.5'' as shown
\item All the positive values above the threshold will be ``True Positives''
\item Negative values above the threshold will be ``False Positives'' as they
are predicted incorrectly as positives.
\item All the negative values below the threshold will be ``True Negatives''
\item Positive values below the threshold will be ``False Negative'' as they
are predicted incorrectly as negatives.
\end{itemize}
\end{column}
\begin{column}[T]{0.5\linewidth}

\begin{center}
\includegraphics[width=0.9\linewidth]{roc2}

\includegraphics[width=0.9\linewidth]{roc3}
\end{center}

\tiny{(Ref: Let's learn about AUC ROC Curve! - GreyAtom)}
\end{column}

\end{columns}
\end{frame}

%%%%%%%%%%%%%%%%%%%%%%%%%%%%%%%%%%%%%%%%%%%%%%%%%%%%%%%%%%%%%%%%%%%%%%%%
\begin{frame}[fragile]\frametitle{What is Sensitivity and Specificity?}

\begin{columns}
\begin{column}[T]{0.5\linewidth}
\begin{itemize}
\item Sensitivity $= \frac{TP}{TP + FN}$: of the actual Yes cases, how many did we catch?
\item Specificity $= \frac{TN}{TN + FP}$: of the actual No cases, how many did we clear?
\item Sensitivity, Recall and TPR are three names for the same quantity.
\item Lowering the threshold raises sensitivity but lowers specificity, and vice versa. A trade-off.
\item We plot (1 - Specificity) instead, as in the graph below.
\end{itemize}
\end{column}
\begin{column}[T]{0.5\linewidth}

\begin{center}
\includegraphics[width=0.7\linewidth]{roc4}

\includegraphics[width=0.7\linewidth]{roc5}
\end{center}

\tiny{(Ref: Let's learn about AUC ROC Curve! - GreyAtom)}
\end{column}

\end{columns}
\end{frame}

%%%%%%%%%%%%%%%%%%%%%%%%%%%%%%%%%%%%%%%%%%%%%%%%%%%%%%%%%%%%%%%%%%%%%%%%
\begin{frame}[fragile]\frametitle{Why do we use (1- Specificity)?}

\begin{columns}
\begin{column}[T]{0.5\linewidth}
\begin{itemize}
\item Let's derive what exactly is (1-Specificity)
\item  Specificity gives us the True Negative Rate
\item  (1 - Specificity) gives us the  False Positive Rate (FPR)
\item So now we are just looking at the positives. 
\item As we increase the threshold,
we decrease the TPR as well as the FPR and when we decrease the
threshold, we are increasing the TPR and FPR.
\end{itemize}
\end{column}
\begin{column}[T]{0.5\linewidth}

\begin{center}
\includegraphics[width=0.8\linewidth]{roc7}
\end{center}

\tiny{(Ref: Let's learn about AUC ROC Curve! - GreyAtom)}
\end{column}

\end{columns}
\end{frame}


%%%%%%%%%%%%%%%%%%%%%%%%%%%%%%%%%%%%%%%%%%%%%%%%%%%%%%%%%%%%%%%%%%%%%%%%
\begin{frame}[fragile]\frametitle{What is AUC?}

\begin{columns}
\begin{column}[T]{0.5\linewidth}
\begin{itemize}
\item The AUC is the area under the ROC curve. 
\item It gives us a good idea
of how well the model performances.
\item AUC ROC indicates how well the probabilities from the positive
classes are separated from the negative classes.
\end{itemize}
\end{column}
\begin{column}[T]{0.5\linewidth}

\begin{center}
\includegraphics[width=0.8\linewidth]{roc6}
\end{center}

\tiny{(Ref: Let's learn about AUC ROC Curve! - GreyAtom)}
\end{column}

\end{columns}
\end{frame}



%%%%%%%%%%%%%%%%%%%%%%%%%%%%%%%%%%%%%%%%%%%%%%%%%%%%%%%%%%%%%%%%%%%%%%%%%
%\begin{frame}[fragile]\frametitle{``Default.csv'' Dataset}
%\begin{itemize}
%\item 10,000 observations
%
%\item 4 variables
%	\begin{itemize}
%	\item Default: ${Yes, No}$ - whether customer defaulted on their debt
%	\item Student: ${Yes, No}$ - whether customer is a student
%	\item Balance: average CC balance
%	\item Income: customer income
%	\end{itemize}
%
%\item Considering Balance as feature and Default as outcome, $\beta_1 = 0.0047$
%\end{itemize}
%%\begin{lstlisting}
%%reg = linear_model.LogisticRegression()
%%reg.fit(default['balance'].reshape(-1,1), default['default'])
%%Coefficients: [[ 0.00478248]]
%%Intercept: [-9.46506555] 
%%\end{lstlisting}
%A one-unit increase in balance is associated with an increase in the log odds of default by 0.0047 units.
%\end{frame}
%
%%%%%%%%%%%%%%%%%%%%%%%%%%%%%%%%%%%%%%%%%%%%%%%%%%%%%%%%%%%%%%%%%%%%%%%%%
%\begin{frame}[fragile]\frametitle{Making Predictions}
%For different Balance(s) what are the chances of Default?
%\begin{center}
%\includegraphics[width=\linewidth]{predlog}
%\end{center}
%Find : Do students have a higher chance of default?
%\end{frame}

% %%%%%%%%%%%%%%%%%%%%%%%%%%%%%%%%%%%%%%%%%%%%%%%%%%%%%%%%%%%%%%%%%%%%%%%%%%%%%%%%%%
% \begin{frame}[fragile]\frametitle{}
% \begin{center}
% {\Large Evaluation Metrics: ROC AUC}
% \end{center}
% \end{frame}

%%%%%%%%%%%%%%%%%%%%%%%%%%%%%%%%%%%%%%%%%%%%%%%%%%%%%%%%%%%%%%%%%%%%%%%%
\begin{frame}[fragile]\frametitle{ROC AUC: A Worked Picture}

\begin{columns}
\begin{column}[T]{0.5\linewidth}
\begin{itemize}
\item ROC : Receiver Operating Characteristics Curve
\item AUC: Area under curve
\item Often used to evaluate performance of a binary classifier
\item Two possible outcomes: Positive and Negative
\item This is histogram (continuous) where red pixels are positive (got admission to IIT!!) and blue pixels are negative (did not get admitted)
\item 250 students (pixels here, for example) got admitted and 250 did not get admitted.
\end{itemize}
\end{column}
\begin{column}[T]{0.5\linewidth}

\begin{center}
\includegraphics[width=0.9\linewidth]{binclass1}
\end{center}

\tiny{(Ref: ROC Curves and Area Under the Curve (AUC) Explained - Data School)}
\end{column}

\end{columns}
\end{frame}

%%%%%%%%%%%%%%%%%%%%%%%%%%%%%%%%%%%%%%%%%%%%%%%%%%%%%%%%%%%%%%%%%%%%%%%%
\begin{frame}[fragile]\frametitle{ROC AUC: Reading the Histogram}

\begin{columns}
\begin{column}[T]{0.5\linewidth}
\begin{itemize}
\item This is result of predictions done on validation set. For which, you have actual values as well.
\item The classifier is like Logistic Regression where you get probabilities as well.
\item Predicted probabilities are seen on the x axis. Y axis is count of students with ACTUAL red(admitted) or blue (not admitted) values
\item E.g. there were 10 students for which you predicted admission probability of 0.1, they did not get admitted shown as negative (blue)
\end{itemize}
\end{column}
\begin{column}[T]{0.5\linewidth}

\begin{center}
\includegraphics[width=0.9\linewidth]{binclass2}
\end{center}

\tiny{(Ref: ROC Curves and Area Under the Curve (AUC) Explained - Data School)}
\end{column}

\end{columns}
\end{frame}

%%%%%%%%%%%%%%%%%%%%%%%%%%%%%%%%%%%%%%%%%%%%%%%%%%%%%%%%%%%%%%%%%%%%%%%%
\begin{frame}[fragile]\frametitle{ROC AUC: A Well-Separated Classifier}

\begin{columns}
\begin{column}[T]{0.5\linewidth}
\begin{itemize}
\item E.g. there were 50 students for which you predicted admission probability of 0.3, they did not get admitted, ie negative (blue) (not shown here, but height of the bell curve is 50 )
\item For 20 students where you predicted 0.5 probability of admission, 10 did not get admitted (blue) but 10 got admitted (red)
\item You can say that the classifier is doing well as it is able to partition/separate two classes reasonably well (only small overlap or ambiguity in the middle)
\end{itemize}
\end{column}
\begin{column}[T]{0.5\linewidth}

\begin{center}
\includegraphics[width=0.9\linewidth]{binclass3}
\end{center}

\tiny{(Ref: ROC Curves and Area Under the Curve (AUC) Explained - Data School)}
\end{column}

\end{columns}
\end{frame}

%%%%%%%%%%%%%%%%%%%%%%%%%%%%%%%%%%%%%%%%%%%%%%%%%%%%%%%%%%%%%%%%%%%%%%%%
\begin{frame}[fragile]\frametitle{ROC AUC: Picking a Threshold at 0.5}

\begin{columns}
\begin{column}[T]{0.5\linewidth}
\begin{itemize}
\item So. we can put threshold (or the partition) at 0.5, anything above that is admitted, and below is not admitted.
\item With this threshold the accuracy looks to be in the range of 90\%.
\item Please note that the ROC curve in the top left is almost touching the left and top boundaries. Its bit away at the corner due to some errors in the middle region.
\end{itemize}
\end{column}
\begin{column}[T]{0.5\linewidth}

\begin{center}
\includegraphics[width=0.9\linewidth]{binclass4}
\end{center}

\tiny{(Ref: ROC Curves and Area Under the Curve (AUC) Explained - Data School)}
\end{column}

\end{columns}
\end{frame}

%%%%%%%%%%%%%%%%%%%%%%%%%%%%%%%%%%%%%%%%%%%%%%%%%%%%%%%%%%%%%%%%%%%%%%%%
\begin{frame}[fragile]\frametitle{ROC AUC: A Poor Classifier}

\begin{columns}
\begin{column}[T]{0.5\linewidth}
\begin{itemize}
\item If the classifier was not good, there would be far more overlap in the middle,ie many misclassifications.
\item Regardless of where you put threshold, your classification accuracy will be bad
\item Note that the ROC curve has shifted far away from the top left corner, so many errors.
\item ROC curve is a plot of True Positive Rate (TPR) on y axis and False Positive Rate (FPR) on x axis for every possible classification threshold (say, 0 to 1)
\end{itemize}
\end{column}
\begin{column}[T]{0.5\linewidth}

\begin{center}
\includegraphics[width=0.9\linewidth]{binclass5}
\end{center}

\tiny{(Ref: ROC Curves and Area Under the Curve (AUC) Explained - Data School)}
\end{column}

\end{columns}
\end{frame}

%%%%%%%%%%%%%%%%%%%%%%%%%%%%%%%%%%%%%%%%%%%%%%%%%%%%%%%%%%%%%%%%%%%%%%%%
\begin{frame}[fragile]\frametitle{ROC AUC: Computing TPR and FPR at 0.8}

\begin{columns}
\begin{column}[T]{0.5\linewidth}
\begin{itemize}
\item TPR : True Positives Predicted divided by All actual positives
\item FPR : False Positives Predicted divided by All actual negatives
\item Both range from 0 to 1.
\item For threshold 0.8, TPR is pixels on right (approx 50) divided by all reds (actual positives) ie 250, so $50/250 = 0.2$,
\item FPR is blue pixels on right of the line (0) divided by all blues ie 250, so $0/250 = 0$
\item Plot point at x = 0, and y = 0.2
\end{itemize}
\end{column}
\begin{column}[T]{0.5\linewidth}

\begin{center}
\includegraphics[width=0.9\linewidth]{binclass6}
\end{center}

\tiny{(Ref: ROC Curves and Area Under the Curve (AUC) Explained - Data School)}
\end{column}

\end{columns}
\end{frame}


%%%%%%%%%%%%%%%%%%%%%%%%%%%%%%%%%%%%%%%%%%%%%%%%%%%%%%%%%%%%%%%%%%%%%%%%
\begin{frame}[fragile]\frametitle{ROC AUC: Computing TPR and FPR at 0.5}

\begin{columns}
\begin{column}[T]{0.5\linewidth}
\begin{itemize}
\item For threshold 0.5, TPR is pixels on right (approx 230) divided by all reds (actual positives) ie 250, so $230/250 = 0.94$,
\item FPR is blue pixels on right of the line (125) divided by all blues ie 250, so $125/250 = 0.5$
\item Plot point at x = 0.5, and y = 0.94
\end{itemize}
\end{column}
\begin{column}[T]{0.5\linewidth}

\begin{center}
\includegraphics[width=0.9\linewidth]{binclass7}
\end{center}

\tiny{(Ref: ROC Curves and Area Under the Curve (AUC) Explained - Data School)}
\end{column}

\end{columns}
\end{frame}

%%%%%%%%%%%%%%%%%%%%%%%%%%%%%%%%%%%%%%%%%%%%%%%%%%%%%%%%%%%%%%%%%%%%%%%%
\begin{frame}[fragile]\frametitle{ROC AUC: Sweeping All Thresholds}

\begin{columns}
\begin{column}[T]{0.5\linewidth}
\begin{itemize}
\item Similarly plot for all possible thresholds. 
\item Thus ROC curve visualizes all possible thresholds.
\item It shows how the classifier is AWAY from the ideal, top left corner
\item The worst classier would be the diagonal straight line.  50\% line ie random chance.
\item So, to distinguish good classifier from a bad one, Area Under Curve is used.
\end{itemize}
\end{column}
\begin{column}[T]{0.5\linewidth}

\begin{center}
\includegraphics[width=0.9\linewidth]{binclass8}
\end{center}

\tiny{(Ref: ROC Curves and Area Under the Curve (AUC) Explained - Data School)}
\end{column}

\end{columns}
\end{frame}

%%%%%%%%%%%%%%%%%%%%%%%%%%%%%%%%%%%%%%%%%%%%%%%%%%%%%%%%%%%%%%%%%%%%%%%%
\begin{frame}[fragile]\frametitle{ROC AUC: Reading the Area}

\begin{columns}
\begin{column}[T]{0.5\linewidth}
\begin{itemize}
\item AUC is Area under ROC curve divided by the whole square, say 0.8 for this example.
\item A poor classifier (diagonal line) will have area 0.5
\item Best classifier will be (corner legs) will have area 1.
\item Class labels are balanced here, so both bell curves have the same height.
\item That balance does not affect the ROC AUC calculation anyway.
\item Choosing threshold is a business decision.
\end{itemize}
\end{column}
\begin{column}[T]{0.5\linewidth}

\begin{center}
\includegraphics[width=0.9\linewidth]{binclass9}
\end{center}

\tiny{(Ref: ROC Curves and Area Under the Curve (AUC) Explained - Data School)}
\end{column}

\end{columns}
\end{frame}

%%%%%%%%%%%%%%%%%%%%%%%%%%%%%%%%%%%%%%%%%%%%%%%%%%%%%%%%%%%%%%%%%%%%%%%%
\begin{frame}[fragile]\frametitle{ROC AUC}

\begin{columns}
\begin{column}[T]{0.5\linewidth}
\begin{itemize}
\item AUC depends only on the \emph{order} of the predicted probabilities.
\item It ignores how confident those probabilities are.
\item So even though model 1 looks far better than model 2, both rank the examples the same way. Same ROC curve, same AUC.
\item AUC alone therefore cannot separate these two models.
\item A better criteria for comparison would be log-loss, which does penalise weak confidence.
\end{itemize}
\end{column}
\begin{column}[T]{0.5\linewidth}

\begin{center}
\includegraphics[width=\linewidth]{logreg19}
\end{center}

{\tiny (Ref: Introduction to Data Science - Analytics Vidhya)}
\end{column}

\end{columns}
\end{frame}

%%%%%%%%%%%%%%%%%%%%%%%%%%%%%%%%%%%%%%%%%%%%%%%%%%%%%%%%%%%%%%%%%%%%%%%%
\begin{frame}[fragile]\frametitle{Log Loss}

\begin{columns}
\begin{column}[T]{0.5\linewidth}
\begin{itemize}
\item As predictions are for output ``1'', probabilities are adjusted first.
\item Then the diff (loss) is computed. Its log is taken.
\item As log is negative, just to make it positive -ve sign is added.
\item This procedure can be concisely put in a formula, as shown
\end{itemize}

\begin{block}{Intuition}
Log Loss is the penalty for being confidently wrong. Predicting 0.9 for a case that turns out to be 0 costs far more than predicting 0.6 for it, because $-\ln$ of a small number blows up. This is why it can separate two models that AUC cannot.
\end{block}
\end{column}
\begin{column}[T]{0.5\linewidth}

\begin{center}
\includegraphics[width=\linewidth]{logreg20}

\includegraphics[width=\linewidth]{logreg21}

\end{center}

{\tiny (Ref: Introduction to Data Science - Analytics Vidhya)}
\end{column}

\end{columns}
\end{frame}

% %%%%%%%%%%%%%%%%%%%%%%%%%%%%%%%%%%%%%%%%%%%%%%%%%%%%%%%%%%%%%%%%%%%%%%%%
% \begin{frame}[fragile]\frametitle{ROC AUC Example}
% Lets look at one example:

% \begin{itemize}

% \item Form a table with y as actual values and $P_+$ predictions probabilities as given by logistic regression.
% \item So, for 1s we will get close to 1 probabilities and for 0s, close to 0.
% \item 3rd example shown is for error.
% \item Sort on Probabilities, like ranking or scoring.
% \end{itemize}

% \begin{center}
% \includegraphics[width=0.35\linewidth]{cf1}
% \end{center}
% \end{frame}

% %%%%%%%%%%%%%%%%%%%%%%%%%%%%%%%%%%%%%%%%%%%%%%%%%%%%%%%%%%%%%%%%%%%%%%%%
% \begin{frame}[fragile]\frametitle{ROC AUC Example}

% \begin{itemize}
% \item Decide a threshold and partition the sorted column 
% \item New column for Predictions where value above partition are 1s and below, 0s.
% \item Decide following terms:
% \begin{itemize}
% \item True Positive Rate (TPR or Recall) $\frac{TP}{TP+FN}$ (shown magenta)
% \item False Positive Rate (FPR) $\frac{FP}{FP+TN}$ (shown red)
% \end{itemize}
% \end{itemize}

% \begin{center}
% \includegraphics[width=0.8\linewidth]{cf2}
% \end{center}
% \end{frame}

% %%%%%%%%%%%%%%%%%%%%%%%%%%%%%%%%%%%%%%%%%%%%%%%%%%%%%%%%%%%%%%%%%%%%%%%%
% \begin{frame}[fragile]\frametitle{ROC AUC Example}

% \begin{itemize}
% \item Vary threshold to get different TPR and FPR
% \item Plot the points with FPR on x and TPR on y, for different thresholds.
% \item Ideal would be with TPR = 1 and FPR = 0. The y axis point. So the whole rectangle is the full area.
% \item Practical is of course less than ideal, as seen. $\frac{8}{9}$
% \item AUC as 0.5 is like random classification. The baseline where TPR and FPR are equal.
% \end{itemize}

% \begin{center}
% \includegraphics[width=0.65\linewidth]{cf3}
% \end{center}
% \end{frame}

% %%%%%%%%%%%%%%%%%%%%%%%%%%%%%%%%%%%%%%%%%%%%%%%%%%%%%%%%%%%%%%%%%%%%%%%%
% \begin{frame}[fragile]\frametitle{ROC AUC Example}

% % \begin{center}
% % \includegraphics[width=0.3\linewidth]{logistic4}
% % \end{center}

% \begin{itemize}
% \item AUC interpretation: If you pick random positive example and a negative one, whats the probability that my model will give higher score to positive (meaning correct class identification) than to the negative.
% \item That value is the area under curve. If its 0.9, thats the probability of getting correct class identification.
% \item Intuition: gives an aggregate measure of performance aggregated across all possible classification thresholds.
% % \item Not IDEAL measure. If errors are more, it calculates misclassification as ROC AUC and thats wrong. (??)
% \end{itemize}
% \end{frame}


% %%%%%%%%%%%%%%%%%%%%%%%%%%%%%%%%%%%%%%%%%%%%%%%%%%%%%%%%%%
% \begin{frame}[fragile]\frametitle{ROC AUC (Recap)}
% \begin{center}
% \includegraphics[width=0.9\linewidth,keepaspectratio]{confmat12}
% \end{center}
% \tiny{(Reference: Introduction to Machine Learning - Andrew Ng)}
% \end{frame}


%%%%%%%%%%%%%%%%%%%%%%%%%%%%%%%%%%%%%%%%%%%%%%%%%%%%%%%%%%%%%%%%%%%%%%%%%%%%%%%%%%
\begin{frame}[fragile]\frametitle{Quick Check: Evaluation Metrics}
\begin{block}{Think About It}
The 99\% accurate ``always predict No Cancer'' model from earlier is useless. What is its AUC, and what does that tell you about which metric to trust on an imbalanced dataset?
\end{block}
\end{frame}

%%%%%%%%%%%%%%%%%%%%%%%%%%%%%%%%%%%%%%%%%%%%%%%%%%%%%%%%%%%%%%%%%%%%%%%%%%%%%%%%%%
\begin{frame}[fragile]\frametitle{Quick Check: Evaluation Metrics}
\begin{block}{Answer}
Its AUC is 0.5, the diagonal, the same as random guessing. That model hands every patient the identical score, so it produces no ordering at all, and AUC measures exactly that ordering. Accuracy was flattered by the 99 to 1 class split; AUC was not, because it sweeps every threshold instead of trusting one. On an imbalanced problem, an accuracy near the majority class share tells you almost nothing, so read the confusion matrix and the AUC before believing a model is good.
\end{block}
\end{frame}
